\documentclass{article}
\usepackage[margin=1.5cm]{geometry}
\usepackage{amsmath}
\usepackage{enumitem}

\setlist{nosep,leftmargin=*}

\begin{document}
\twocolumn

\title{Chapters 1--5 Warm-Up: Python Review}
\author{Prof. Jordan C. Hanson}
\date{}

\maketitle

\section{Expressions and Variables}

\begin{enumerate}
\item Predict the output without running the code.  Then check your
answer in Python.
\begin{verbatim}
x = 8
y = 3
print(x + y)
print(x * y)
print(x / y)
print(x % y)
\end{verbatim}
\item Write a Python expression that converts a length provided by the user in inches to meters.
\end{enumerate}

\section{Strings and Input}

\begin{enumerate}
\item Create a program that asks the user for a name and an age, and prints the results in the following format using an f-string: ``name: age.''
\end{enumerate}

\section{Lists and Tuples}

\begin{enumerate}
\item Consider
\begin{verbatim}
student = ("Maya", 20681001, 4)
course = ["Calculus",["Maya", "Liam", "Ava"]]
\end{verbatim}
What is produced by each expression?
\begin{verbatim}
student[1]
student[2]
course[0]
course[1][2]
\end{verbatim}
\item Add \texttt{"Noah"} to the student list inside \texttt{course}.
\end{enumerate}

\section{Boolean Logic and Decisions}

\begin{enumerate}
\item Let
\begin{verbatim}
score = 3
submitted = True
\end{verbatim}
Predict the value of each expression:
\begin{verbatim}
score >= 4
score >= 2 and submitted
score == 0 or score == 1
not submitted
\end{verbatim}
\item Complete the decision:
\begin{verbatim}
if score >= 4:
    print("Calculus")
elif __________:
    print("Trigonometry")
else:
    print("Algebra")
\end{verbatim}
\end{enumerate}

\section{Loops}

\begin{enumerate}
\item What values are printed?
\begin{verbatim}
for i in range(4):
    print(i)
\end{verbatim}
How many times does the loop execute?
\item Complete the following loop so that it prints integers in order:
\begin{verbatim}
i = 0
while i < 4:
    print(i)
    __________
\end{verbatim}
\item Consider
\begin{verbatim}
scores = [5, 1, 3, 4, 2]
\end{verbatim}
Write a \texttt{for} loop that prints \texttt{"Calculus"} whenever the score is 4 or 5.
\end{enumerate}

\section{Challenge}

\begin{enumerate}

\item Define the following \textit{diagonal} matrix:

\begin{verbatim}
A = [
    [1, 0, 0],
    [0, 1, 0],
    [0, 0, 1]
]
\end{verbatim}

Suppose we did not know the matrix was diagonal.  Create nested \verb+for+ loops and conditional logic that demonstrate the matrix is diagonal.  For example:

\begin{verbatim}
is_diagonal = True
for row in range(2):
    for column in range(3):
        if(...):
        		is_diagonal = False
        	...
\end{verbatim}
\end{enumerate}

\end{document}